\section{Smooth summaries at a spatial scale}
\label{sec:channels}

The connection forms will be built from smooth summaries of probability
distributions on the group.  A summary must retain information about spatial
concentration while changing in a controlled way when mass moves between two
group elements.  We construct two versions.  The first allows the spatial
radius to increase by a factor of six and has a derivative bound independent
of that radius.  The second retains the radius and has a polynomial dependence
on it.  Their different bounds will be used at different stages of the
averaging construction.

Throughout this section, $G$ is a finitely generated group with a fixed left
invariant word metric $d$.  Write $B(y,s)=\{h:d(h,y)\le s\}$.  Counting words
in a finite symmetric generating set gives a constant $c_d$ such that
\begin{equation}
 |B(y,s)|\le \exp(c_d s)\qquad(s\ge1).
 \label{ch:ball-growth}
\end{equation}
Indeed, if the generating set has $m$ elements, the ball is covered by words
of length at most $\lfloor s\rfloor$, and their number is at most
$(m+1)^{\lfloor s\rfloor+1}$.  Increasing $c_d$ also covers the trivial
group.  Constants denoted by $C$ may depend on this word metric but not on a
radius, an accuracy parameter, or the size of an input support.

Let $\mathcal P_f(G)$ be the set of finitely supported probability
distributions on $G$.  For $r\ge1$ set
\begin{equation}
 \begin{split}
 \chi_r(s)&=\min\{1,\max\{0,2-s/r\}\},\qquad s\ge0,\\
 E_r(P,y)&=\sum_{h\in G}P_h\chi_r(d(h,y)),\\
 q_r(P)&=\max_{y\in G}E_r(P,y),\qquad
 \phi_r(P)=(q_r(P)-a_*)_+,\qquad a_*=0.80.
 \end{split}
 \label{ch:concentration}
\end{equation}
The maximum exists: the score vanishes outside the finite $2r$ neighborhood
of $\supp P$.  It lies in $[0,1]$, so $0\le\phi_r\le0.20$.  The functions
$q_r$ and $\phi_r$ need not be smooth; they will only be used in inequalities.

\begin{definition}
\label{ch:channel-definition}
A \emph{channel} with label set $\Lambda$, on which $G$ acts, assigns to
each $P\in\mathcal P_f(G)$ a finitely supported subprobability
$Y(P)=(Y(P)_\omega)_{\omega\in\Lambda}$.  Thus $Y(P)_\omega\ge0$ and
$\sum_\omega Y(P)_\omega\le1$.  Translation of weights means
$(sP)_h=P_{s^{-1}h}$ and $(sY)_\omega=Y_{s^{-1}\omega}$.
We require left equivariance,
$Y(sP)=sY(P)$ for $s\in G$, and the following smoothness property.  For each
finite $S\subset G$, the restriction to probabilities supported in $S$ has
its labels in a fixed finite subset of $\Lambda$, and its coordinate
functions extend smoothly to a neighborhood of that probability simplex in
$\R^S$.
\end{definition}

In a transfer direction $\dot P=\delta_h-\delta_{h'}$, derivatives are
computed after enlarging $S$ to contain $h,h'$.  The channels we construct
will have compatible formulas on these enlarged spaces, and their
derivatives will vanish at zero output coordinates.  For such a channel
define its Fisher seminorm in
this direction by
\begin{equation}
 \mathcal F_Y(P;h,h')
 =\left(\sum_{\omega:Y(P)_\omega>0}
       \frac{|\dot Y(P)_\omega|^2}{Y(P)_\omega}\right)^{1/2}.
 \label{ch:fisher-definition}
\end{equation}
The zero-coordinate assertion is part of the construction, and will be
proved; it is not a convention for discarding a nonzero derivative.

If $A$ has label set $G$ and $L$ has label set $\Lambda$, a \emph{cross
kernel} is a function $\kappa:G\times\Lambda\to[0,1]$ satisfying
$\kappa(sy,s\omega)=\kappa(y,\omega)$.  It is fixed independently of $P$.
It gives the predictions
\begin{equation}
 \operatorname{pred}(P,y)
   =\sum_{\omega\in\Lambda}\kappa(y,\omega)L(P)_\omega.
 \label{ch:prediction-definition}
\end{equation}
Exact predictions alone would be easy: $L(P)=P$ and
$\kappa(y,h)=\chi_r(d(h,y))$ give $\operatorname{pred}(P,y)=E_r(P,y)$.
When $G$ has at least two elements, this choice has no uniform Fisher bound:
a transfer between distinct elements, adding mass to an entry of size
$t>0$, contributes $1/t$ to the squared seminorm, which diverges as
$t\downarrow0$.  The approximation allowed in the next proposition permits
uniform derivative control, including at probability faces.

\begin{proposition}[Spatial summary channels]
\label{ch:summary}
Fix $r\ge1$ and $0<\epsilon\le10^{-3}$.  There are channels $A,L$ and a
cross kernel $\kappa$ with either of the following two choices of radius:
the early version has $r'=6r$, and the later version has $r'=r$.
In both versions $A$ is labelled by $G$, and for every
$P\in\mathcal P_f(G)$ and $y\in G$,
\begin{align}
 \sum_{z\in G}A(P)_z\bigl(E_r(P,z)-a_*\bigr)
       &\ge\phi_r(P)-10\epsilon,
       \label{ch:selection-bound}\\
 \operatorname{pred}(P,y)
       &\le E_{r'}(P,y)+10\epsilon,
       \label{ch:prediction-upper}\\
 \operatorname{pred}(P,y)
       &\ge E_r(P,y)-10\epsilon
       \qquad\text{if }A(P)_y>0.
       \label{ch:prediction-lower}
\end{align}
For $Y=A$ or $Y=L$, derivatives in every transfer direction vanish at zero
coordinates and satisfy
\begin{equation}
 \mathcal F_Y(P;h,h')\le C(1+d(h,h'))
 \begin{cases}
   \epsilon^{-10},&\text{early version},\\
   r^{10}\epsilon^{-10},&\text{later version}.
 \end{cases}
 \label{ch:summary-fisher}
\end{equation}
Here $C$ depends only on the fixed word metric.  The numerical error constant
$10$ in \eqref{ch:selection-bound}--\eqref{ch:prediction-lower} is independent
of the group and of the parameters.
\end{proposition}

To see why these guarantees are suited to averaging, assume them for the
moment.  Take a finite family $P^{(\nu)}$ of input probabilities and weights
$\theta_\nu\ge0$ with $\sum_\nu\theta_\nu=1$, and write
\[
 \overline P=\sum_\nu\theta_\nu P^{(\nu)},\qquad
 \overline A_y=\sum_\nu\theta_\nu A(P^{(\nu)})_y,\qquad
 \overline{\operatorname{pred}}_y
       =\sum_\nu\theta_\nu\operatorname{pred}(P^{(\nu)},y).
\]
The selection and lower prediction bounds give
$\sum_yA(P)_y(\operatorname{pred}(P,y)-a_*)\ge\phi_r(P)-20\epsilon$.
On the other hand, the upper prediction bound and linearity of $E_{r'}$
give
$\overline{\operatorname{pred}}_y
\le E_{r'}(\overline P,y)+10\epsilon$.
Since $\overline A$ is a subprobability, its pairing with
$E_{r'}(\overline P,\cdot)-a_*$ is at most
$\phi_{r'}(\overline P)$; if the maximal centered score is negative, that
pairing is nonpositive.  Finally,
$\sum_\nu\theta_\nu\sum_yA(P^{(\nu)})_y=\sum_y\overline A_y$, so
the terms containing $a_*$ cancel and
\begin{equation}
 \begin{split}
 &\sum_\nu\theta_\nu\sum_y
       A(P^{(\nu)})_y\operatorname{pred}(P^{(\nu)},y)
       -\sum_y\overline A_y\overline{\operatorname{pred}}_y\\
 &\hspace{2em}\ge
       \sum_\nu\theta_\nu\phi_r(P^{(\nu)})
       -\phi_{r'}(\overline P)-30\epsilon.
 \end{split}
 \label{ch:averaging-drop}
\end{equation}
Thus loss of concentration under averaging gives a lower bound for the
difference of the two pairings.  The same calculation applies to time
averages of fields with a common finite support, which is how these
estimates will enter the connection construction.

We first construct $A$ and a second distribution $D$ on possible centers.
The channel $A$ concentrates near a maximizing center when the maximal score
exceeds $a_*$.  The thresholds for $D$ are lower, so that $D$ already has
almost full mass whenever $A$ is nonzero.  The two versions of $L$ will then
use this same $D$ to predict scores at the selected centers.

\subsection{Smooth selection of centers}

The number of centers that can carry substantial score has a bound
independent of $|\supp P|$.

\begin{lemma}
\label{ch:center-geometry}
For $P\in\mathcal P_f(G)$, the set
$\{y:E_r(P,y)>1/2\}$ has diameter at most $4r$ and cardinality at most
$\exp(\gamma r)$, where $\gamma$ depends only on the word metric.
\end{lemma}

\begin{proof}
The inequality $E_r(P,y)>1/2$ implies
$P(\{h:d(h,y)<2r\})>1/2$.  Any two such sets intersect, since their
masses add to more than one.  Their centers consequently have distance less
than $4r$.  If the set of centers is nonempty, choose one of them.  All the
others lie in its $4r$ ball, and \eqref{ch:ball-growth} gives the claimed
bound, for example with $\gamma=4c_d$.  The empty case is immediate.
\end{proof}

Choose smooth functions $\psi_A,\psi_D:\R\to[0,1]$ such that
\begin{equation}
 \begin{array}{c|cc}
       &\text{identically zero on}&\text{identically one on}\\ \hline
 \psi_A&(-\infty,0.70]&[0.75,\infty)\\
 \psi_D&(-\infty,0.58]&[0.62,\infty).
 \end{array}
 \label{ch:cutoffs}
\end{equation}
For completeness, a smooth increasing cutoff from zero to one on $[0,1]$
is $\eta(t)/(\eta(t)+\eta(1-t))$, where
$\eta(t)=e^{-1/t}$ for $t>0$ and $\eta(t)=0$ for $t\le0$.
Every right derivative of $\eta$ at zero vanishes: away from zero it is
$e^{-1/t}$ times a polynomial in $1/t$, which tends to zero as $t\downarrow0$.
The denominator in the displayed quotient is everywhere positive, and
affine changes of variable give \eqref{ch:cutoffs}.  Their first derivatives
are bounded by a fixed numerical constant, since they vanish outside fixed
compact transition intervals.

Take $N=M_0r\epsilon^{-2}$ and define
\begin{equation}
 \begin{split}
 w^A_y&=\psi_A(E_r(P,y))^2
              e^{N(E_r(P,y)-0.80)},\\
 w^D_y&=\psi_D(E_r(P,y))^2
              e^{N(E_r(P,y)-0.66)},\\
 A(P)_y&=\frac{w^A_y}{1+\sum_z w^A_z},\qquad
 D(P)_y=\frac{w^D_y}{1+\sum_z w^D_z}.
 \end{split}
 \label{ch:selector-formulas}
\end{equation}
Here $M_0$ is fixed in terms of the metric, large enough that for all allowed
$r,\epsilon$,
\begin{equation}
 (1+e^{\gamma r})e^{-N\epsilon}\le\epsilon,
 \qquad e^{\gamma r-0.04N}\le1.
 \label{ch:exponential-choices}
\end{equation}
For example, any $M_0\ge25\gamma+4$ suffices.  Indeed,
$\log(1+e^{\gamma r})+\log(1/\epsilon)
\le\gamma r+\log2+1/\epsilon\le(\gamma+2)r/\epsilon$,
and $N\ge25\gamma r$.

The additional $1$ in each denominator may be regarded as the weight of an
unused null label.  It allows the channel to have small total mass when no
center reaches the exponential threshold.  Each displayed sum is finite by
Lemma~\ref{ch:center-geometry}.

\begin{lemma}[Selector estimates]
\label{ch:selector-estimates}
The formulas \eqref{ch:selector-formulas} define smooth equivariant
channels.  They satisfy
\begin{equation}
 \sum_y A(P)_y(E_r(P,y)-a_*)\ge\phi_r(P)-3\epsilon.
 \label{ch:selector-concentration}
\end{equation}
Whenever $q_r(P)>0.70$,
\begin{equation}
 \sum_zD(P)_z\ge1-\epsilon,
 \qquad
 \sum_zD(P)_zE_r(P,z)\ge q_r(P)-2\epsilon.
 \label{ch:D-concentration}
\end{equation}
Both channels have zero derivatives at zero coordinates, and
\begin{equation}
 \mathcal F_A(P;h,h')+\mathcal F_D(P;h,h')
       \le C\epsilon^{-2}d(h,h').
 \label{ch:selector-fisher}
\end{equation}
\end{lemma}

\begin{proof}
Fix $P$ and abbreviate $E_r(P,y)$ to $E_y$ and $q_r(P)$ to $q$.
Suppose first that $q>a_*+\epsilon$.  At a maximizing center the cutoff
$\psi_A$ is one, and its weight is $e^{N(q-a_*)}$.
The normalized mass of centers with $E_y<q-\epsilon$, together with the
null mass, is at most
$(1+e^{\gamma r})e^{-N\epsilon}\le\epsilon$.
On the remaining centers $E_y-a_*\ge q-a_*-\epsilon$; on all centers
$E_y-a_*\ge-1$.  Their weighted sum is therefore at least
$q-a_*-3\epsilon$.  If $q\le a_*+\epsilon$, the total mass at centers
with $E_y<a_*-\epsilon$ is at most $e^{\gamma r}e^{-N\epsilon}\le\epsilon$,
using the $1$ in the denominator.  All other negative contributions are
bounded below by $-\epsilon$ times their mass.  The sum is at least
$-2\epsilon$, whereas $\phi_r(P)\le\epsilon$.  This proves
\eqref{ch:selector-concentration} in both cases.

Now assume $q>0.70$.  The cutoff $\psi_D$ is one at a maximizing center.
The same comparison, with threshold $0.66$, shows that the total mass of
the null label and the centers with $E_z<q-\epsilon$ is at most $\epsilon$.
Indeed the null-weight ratio is $e^{-N(q-0.66)}\le e^{-N\epsilon}$.
Thus the nonnull mass is at least $1-\epsilon$, and the score average is
at least $(1-\epsilon)(q-\epsilon)\ge q-2\epsilon$.  This proves
\eqref{ch:D-concentration}.

We next prove the derivative estimate.  Put $d_0=d(h,h')$.  The function
$\chi_r$ is $1/r$-Lipschitz, so
\begin{equation}
 |\dot E_y|
 =|\chi_r(d(h,y))-\chi_r(d(h',y))|
 \le d_0/r.
 \label{ch:score-derivative}
\end{equation}
The following calculation applies to either cutoff, with its corresponding
threshold $\theta\in\{0.80,0.66\}$.  Set
\[
 f_y=\psi(E_y)e^{N(E_y-\theta)/2},\qquad
 F=(1,(f_y)_y),\qquad Z=\|F\|^2=1+\sum_y f_y^2.
\]
The nonnull coordinates of $F/\|F\|$ are the smooth nonnegative square
roots of the normalized weights.  Differentiating a normalized Euclidean
vector gives
\begin{equation}
 \left\|\frac{d}{dt}\frac{F}{\|F\|}\right\|^2
 \le\frac{\|\dot F\|^2}{Z}.
 \label{ch:normalization-derivative}
\end{equation}
To see this, the derivative is $\|F\|^{-1}$ times the orthogonal
projection of $\dot F$ onto $F^\perp$.

Differentiating $f_y$ and using $(u+v)^2\le2u^2+2v^2$ gives
\begin{equation}
 \frac{\|\dot F\|^2}{Z}
 \le C\frac{d_0^2}{r^2}
 \left(
  \frac{\sum_y|\psi'(E_y)|^2e^{N(E_y-\theta)}}{Z}
   +N^2\frac{\sum_y f_y^2}{Z}
 \right).
 \label{ch:cutoff-derivative-estimate}
\end{equation}
The second fraction is at most one.  A nonzero cutoff derivative requires
$E_y\le0.75$ for $A$ or $E_y\le0.62$ for $D$.  In either case
$E_y-\theta\le-0.04$.  Lemma~\ref{ch:center-geometry} bounds the number
of these terms by $e^{\gamma r}$.  The first fraction is consequently at
most $Ce^{\gamma r-0.04N}\le C$.  Equations
\eqref{ch:normalization-derivative} and
\eqref{ch:cutoff-derivative-estimate} now bound the squared norm of the
square-root derivative by $C(1+N^2)d_0^2/r^2$.  At a positive normalized
weight $Y_y$,
$|\dot Y_y|^2/Y_y=4|\frac{d}{dt}\sqrt{Y_y}|^2$.
Substituting $N=M_0r\epsilon^{-2}$ proves
\eqref{ch:selector-fisher}.

These computations also justify regularity at zero coordinates.  A smooth
nonnegative function has zero derivative wherever it vanishes.  Thus
$\psi(E_y)=0$ implies $\dot f_y=0$, even in a transfer direction that
leaves a probability face.  Hence the normalized weight and its square root
both have zero derivative there.

Finally, on distributions supported in a fixed finite $S$, only centers
within its $2r$ neighborhood can have a nonzero score.  The formulas in
\eqref{ch:selector-formulas} are therefore finite smooth formulas on
$\R^S$, with positive denominators.  Enlarging $S$ adds the same formulas
for the additional coordinates.  This proves the stated smoothness and
compatibility at faces.  Left equivariance follows from the left invariance
of $d$.
\end{proof}

The final averaging stage will only need $D$, together with its concentration
on a bounded set of labels.  We record that consequence separately.

\begin{lemma}[Terminal channel]
\label{ch:terminal}
For each $r\ge1$ and $0<\epsilon\le10^{-3}$, the channel $D$ in
\eqref{ch:selector-formulas} is smooth, left equivariant, and has total
mass at most one.  Its nonempty support has diameter at most $4r$ and
cardinality at most $e^{\gamma r}$, with $\gamma$ depending only on the
word metric.  If $q_r(P)>0.70$, its mass is at least
$1-\epsilon\ge1/2$.  Its transfer derivatives vanish at zero coordinates
and satisfy
\[
 \mathcal F_D(P;h,h')\le C\epsilon^{-2}d(h,h')
       \le C\epsilon^{-10}(1+d(h,h')).
\]
In particular, if $\phi_r(P)>0$, then
\begin{equation}
 \sum_zD(P)_z^2\ge\frac14e^{-\gamma r}.
 \label{ch:terminal-square-mass}
\end{equation}
\end{lemma}

\begin{proof}
Every center in the support of $D$ has score greater than $0.58>1/2$.
Apply Lemma~\ref{ch:center-geometry} and
Lemma~\ref{ch:selector-estimates}.  If $\phi_r(P)>0$, then
$q_r(P)>0.80$, so $\sum_zD_z\ge1/2$.  The Cauchy--Schwarz inequality on
the finite support gives
$\sum_zD_z^2\ge(\sum_zD_z)^2/|\supp D|$, proving
\eqref{ch:terminal-square-mass}.
\end{proof}

\subsection{Predictions with an enlarged radius}

The early channel can use the average score of the centers chosen by $D$.
All these centers lie near any center selected by $A$.  Enlarging the radius
allows us to bound each score contributing at $y$ by $E_{6r}(P,y)$.
Define
\begin{equation}
 \Lambda=G,\qquad
 L(P)_z=D(P)_zE_r(P,z),\qquad
 \kappa(y,z)=\mathbf1_{\{d(y,z)\le4r\}}.
 \label{ch:early-channel}
\end{equation}
This is a subprobability channel, since $0\le E_r(P,z)\le1$.
When $d(y,z)\le4r$, every $h$ with $\chi_r(d(h,z))>0$ has
$d(h,y)<6r$.  On this set $\chi_{6r}(d(h,y))=1$, and therefore
\begin{equation}
 E_r(P,z)\le E_{6r}(P,y).
 \label{ch:radius-domination}
\end{equation}
After summing against $D$, this proves
$\operatorname{pred}(P,y)\le E_{6r}(P,y)$.

If $A(P)_y>0$, then $E_r(P,y)>0.70$ and $q_r(P)>0.70$.  For every
$z$ with $D(P)_z>0$, Lemma~\ref{ch:center-geometry} gives
$d(y,z)\le4r$.  Hence every such $z$ contributes to the prediction, and
\eqref{ch:D-concentration} gives
\[
 \operatorname{pred}(P,y)
 =\sum_zD(P)_zE_r(P,z)
 \ge q_r(P)-2\epsilon
 \ge E_r(P,y)-2\epsilon.
\]

To estimate derivatives, write $E_z=E_r(P,z)$ again.  Where $D_z>0$,
$E_z>0.58$.  Differentiating $L_z=D_zE_z$ yields
\begin{equation}
 \sum_{z:L_z>0}\frac{|\dot L_z|^2}{L_z}
 \le 2\sum_{z:D_z>0}E_z\frac{|\dot D_z|^2}{D_z}
       +2\sum_{z:D_z>0}D_z\frac{|\dot E_z|^2}{E_z}
 \le C\epsilon^{-4}d(h,h')^2.
 \label{ch:early-fisher}
\end{equation}
In the second sum only $D_z>0$ is relevant; its denominators are thus bounded
below by $0.58$.  We used \eqref{ch:score-derivative},
\eqref{ch:selector-fisher}, and $r\ge1$.  At a zero coordinate $D_z=0$
and $\dot D_z=0$, so $\dot L_z=0$.  The same finite formulas prove
smoothness at faces.  This establishes all the early-version claims of
Proposition~\ref{ch:summary}.

\subsection{Accurate finite sample strings}

To keep the radius equal to $r$, an average of center scores is insufficient:
we must predict the score separately at each $y$.  A finite string of sample
points approximates all scores in a bounded region simultaneously.  Random
sampling will prove the existence of an accurate string.  We will then put
smooth positive weights on \emph{all} strings in a fixed finite alphabet.
This distinction is useful at probability faces: the defining weights will
not be products of possibly vanishing input probabilities.

Fix a center $z\in G$.  Introduce a null symbol $\partial$, fixed by the
group action, and the alphabet
\[
 \mathcal A_z=B(z,20r)\cup\{\partial\}.
\]
Define a probability $Q^z$ on this alphabet by
\begin{equation}
 P^z_h=P_h\chi_{10r}(d(h,z)),\qquad
 Q^z_h=P^z_h\ (h\ne\partial),\qquad
 Q^z_\partial=1-\sum_hP^z_h.
 \label{ch:truncated-probability}
\end{equation}
The notation $E_r(P^z,y)$ has the same linear meaning as in
\eqref{ch:concentration}, although $P^z$ can have mass less than one.
This truncation provides a fixed finite alphabet while preserving every
score needed for the lower prediction bound.  Indeed, if
$A(P)_yD(P)_z>0$, Lemma~\ref{ch:center-geometry} gives $d(y,z)\le4r$.
Every $h$ contributing to $E_r(P,y)$ then has $d(h,z)<6r$, where the
factor $\chi_{10r}(d(h,z))$ is one.  Consequently
\begin{equation}
 E_r(P^z,y)=E_r(P,y)
 \qquad\text{if }A(P)_yD(P)_z>0.
 \label{ch:truncation-exactness}
\end{equation}
For arbitrary $y$, we still have $E_r(P^z,y)\le E_r(P,y)$, which is the
direction needed for the upper bound.  Thus accurate predictions for $P^z$
will suffice.  The Lipschitz taper in \eqref{ch:truncated-probability} will
also keep their transfer derivatives proportional to $d(h,h')/r$.

For a positive integer $b$ and a string $w=(w_1,\ldots,w_b)\in\mathcal A_z^b$
let
\begin{equation}
 H_y(w)=\frac1b\sum_{i:w_i\ne\partial}\chi_r(d(w_i,y)),
 \qquad
 \Delta(P;z,w)=\max_{y\in B(z,22r)}
          |H_y(w)-E_r(P^z,y)|.
 \label{ch:string-discrepancy}
\end{equation}
Both scores vanish when $d(y,z)>22r$, by the triangle inequality.  Thus
$\Delta$ also controls the discrepancy at every $y\in G$.  Always
$0\le\Delta\le1$.

\begin{lemma}
\label{ch:accurate-string}
There is a constant $C$, depending only on the word metric, and an integer
$b$ depending only on $r,\epsilon$, with
$1\le b\le Cr\epsilon^{-4}$, such that for every $P$ and $z$ some
$w\in\mathcal A_z^b$ satisfies $\Delta(P;z,w)\le\epsilon$.
\end{lemma}

\begin{proof}
Take $b$ independent samples with distribution $Q^z$.  For a fixed test
center $y$, each summand in $H_y$ lies in $[0,1]$ and has mean
$E_r(P^z,y)$.  The bounded-variable concentration estimate of
Hoeffding~\cite{Hoeffding1963} gives the needed probability bound.  We
include a direct exponential-moment proof, with sufficient constants.
If a real random variable $X$ has mean zero and $|X|\le1$, then, for
$|t|\le1$,
\[
 \mathbf E e^{tX}
 \le1+\sum_{k\ge2}\frac{|t|^k}{k!}
 \le1+t^2\le e^{t^2}.
\]
For independent copies $X_1,\ldots,X_b$, Markov's inequality and
independence, with $t=\epsilon/2$, imply
\[
 \mathbf P\!\left(\frac1b\sum_iX_i>\epsilon\right)
 \le e^{-tb\epsilon}\prod_i\mathbf E e^{tX_i}
 \le e^{-b\epsilon^2/4}.
\]
Applying the same argument to $-X_i$ gives the bound
$2e^{-b\epsilon^2/4}$ for absolute error exceeding $\epsilon$.

By \eqref{ch:ball-growth}, there is a fixed $c_1$ such that the number of
test centers is at most $e^{c_1r}$.  A union bound therefore gives
\[
 \mathbf P\bigl(\Delta(P;z,w)>\epsilon\bigr)
 \le2e^{c_1r-b\epsilon^2/4}.
\]
Choose $b=\lceil M_1r\epsilon^{-4}\rceil$ with a sufficiently large
fixed $M_1$.  The right side is strictly less than one for all permitted
$r,\epsilon$.  At least one string is accurate.  Since
$r\epsilon^{-4}\ge1$, rounding preserves the asserted upper bound on $b$.
\end{proof}

\subsection{Smooth weights on strings}

Fix the integer $b$ supplied by Lemma~\ref{ch:accurate-string}.  To assign
smooth weights, replace the maximum of the signed errors by the logarithm
of a sum of exponentials.  Put
\begin{equation}
 S(P;z,w)=\frac1\lambda\log
 \left(\sum_{y\in B(z,22r)}\sum_{\sigma\in\{-1,1\}}
   \exp\bigl(\lambda\sigma(H_y(w)-E_r(P^z,y))\bigr)\right).
 \label{ch:smooth-discrepancy}
\end{equation}
Choose $\lambda>0$, depending only on $r,\epsilon$, so large that the
logarithm of the number of signed tests is at most $\lambda\epsilon$.
The number of tests is independent of $z$, since left translation identifies
the balls.  The elementary inequalities
$e^{\lambda\max x_i}\le\sum_i e^{\lambda x_i}
\le M e^{\lambda\max x_i}$ give
\begin{equation}
 \Delta(P;z,w)\le S(P;z,w)\le\Delta(P;z,w)+\epsilon.
 \label{ch:smoothing-error}
\end{equation}

The first derivative of this smoothed maximum has a particularly simple
bound.  The test function
$h\mapsto\chi_{10r}(d(h,z))\chi_r(d(h,y))$ is
$11/(10r)$-Lipschitz: both factors have absolute value at most one, and
their Lipschitz constants add.  Thus every signed error in
\eqref{ch:smooth-discrepancy} has transfer derivative of absolute value
at most $11d(h,h')/(10r)$.  Differentiation of the logarithm gives a convex
combination of these derivatives; the factor $\lambda$ cancels.
Consequently
\begin{equation}
 |\dot S(P;z,w)|\le\frac{11}{10r}d(h,h').
 \label{ch:smooth-discrepancy-derivative}
\end{equation}
In particular this estimate does not grow with the number of tests or the
smoothing parameter.

Now define a probability distribution on the entire fixed set of strings by
\begin{equation}
 \pi_z(P)_w
 =\frac{e^{-N_1S(P;z,w)}}{\sum_{v\in\mathcal A_z^b}e^{-N_1S(P;z,v)}},
 \qquad N_1=M_2rb\epsilon^{-2}.
 \label{ch:gibbs-strings}
\end{equation}
Every coordinate is strictly positive.  The constant $M_2$ will be fixed
in terms of the metric.

\begin{lemma}
\label{ch:string-estimates}
For $M_2$ sufficiently large, uniformly in $P,z,y$,
\begin{align}
 \sum_w\pi_z(P)_w\Delta(P;z,w)&\le5\epsilon,
       \label{ch:mean-discrepancy}\\
 \left|\sum_w\pi_z(P)_wH_y(w)-E_r(P^z,y)\right|
       &\le5\epsilon.
       \label{ch:mean-score}
\end{align}
The transfer derivatives satisfy
\begin{equation}
 \left(\sum_w\frac{|\dot\pi_z(P)_w|^2}{\pi_z(P)_w}\right)^{1/2}
       \le C N_1\frac{d(h,h')}{r}.
 \label{ch:string-fisher}
\end{equation}
\end{lemma}

\begin{proof}
There are at most $e^{c_2rb}$ strings, for a constant $c_2$ fixed by the
metric.  By Lemma~\ref{ch:accurate-string} and
\eqref{ch:smoothing-error}, at least one string has $S\le2\epsilon$.
Every string with $\Delta>4\epsilon$ has $S>4\epsilon$, and hence a
weight at most $e^{-2N_1\epsilon}$ times the weight of that particular
string.  Therefore
\[
 \sum_{w:\Delta(P;z,w)>4\epsilon}\pi_z(P)_w
 \le e^{c_2rb-2N_1\epsilon}\le\epsilon
\]
once $M_2$ is sufficiently large.  One fixed choice works for all
$r\ge1$, $b\ge1$, and $\epsilon\le10^{-3}$, as is seen by taking
logarithms and using $\log(1/\epsilon)\le1/\epsilon$.
Since $\Delta\le1$, splitting the expectation at $4\epsilon$ proves
\eqref{ch:mean-discrepancy}.  The discrepancy controls each test center,
and both scores vanish outside the test ball, proving
\eqref{ch:mean-score} for every $y$.

Differentiating the normalized exponential in \eqref{ch:gibbs-strings}
gives
\[
 \frac{\dot\pi_z(P)_w}{\pi_z(P)_w}
 =-N_1\left(\dot S(P;z,w)
       -\sum_v\pi_z(P)_v\dot S(P;z,v)\right).
\]
Equation~\eqref{ch:smooth-discrepancy-derivative} bounds its absolute value
by $C N_1d(h,h')/r$.  Squaring and summing against the probability
$\pi_z(P)$ proves \eqref{ch:string-fisher}.
\end{proof}

\subsection{Predictions at the original radius}

We can now finish the later channel.  Its label set consists of pairs
$(z,w)$, with $z\in G$ and $w\in\mathcal A_z^b$.  Translation acts on
the center and on every nonnull entry of the string.  Define
\begin{equation}
 L(P)_{z,w}=D(P)_z\pi_z(P)_w,
 \qquad \kappa(y,(z,w))=H_y(w).
 \label{ch:later-channel}
\end{equation}
The total mass of $L$ is $\sum_zD(P)_z\le1$, and the kernel takes values
in $[0,1]$.  All formulas commute with left translation.

By \eqref{ch:mean-score}, and since $P^z_h\le P_h$,
\[
 \operatorname{pred}(P,y)
 \le\sum_zD(P)_z\bigl(E_r(P^z,y)+5\epsilon\bigr)
 \le E_r(P,y)+5\epsilon.
\]
Suppose next that $A(P)_y>0$.  For every $z$ with $D(P)_z>0$,
\eqref{ch:truncation-exactness} identifies the predicted truncated score
with $E_r(P,y)$.  Moreover $q_r(P)>0.70$, and
\eqref{ch:D-concentration} gives
$\sum_zD(P)_z\ge1-\epsilon$.  It follows that
\[
 \operatorname{pred}(P,y)
 \ge E_r(P,y)\sum_zD(P)_z-5\epsilon\sum_zD(P)_z
 \ge E_r(P,y)-6\epsilon.
\]
These are the required prediction bounds with $r'=r$.

For the derivative bound it is useful to retain the product structure
$L_{z,w}=D_z\pi_z(w)$.  Since $\sum_w\pi_z(w)=1$ and
$\sum_w\dot\pi_z(w)=0$, expansion of the square gives the exact identity
\begin{equation}
 \sum_{z,w:L_{z,w}>0}\frac{|\dot L_{z,w}|^2}{L_{z,w}}
 =\sum_{z:D_z>0}\frac{|\dot D_z|^2}{D_z}
   +\sum_zD_z\sum_w\frac{|\dot\pi_z(w)|^2}{\pi_z(w)}.
 \label{ch:product-fisher}
\end{equation}
All weights and derivatives here are real; the mixed term vanishes by the
displayed zero-sum identity.  Using \eqref{ch:selector-fisher},
\eqref{ch:string-fisher}, and $b\le Cr\epsilon^{-4}$ yields
\begin{equation}
 \mathcal F_L(P;h,h')
 \le C\bigl(\epsilon^{-2}+N_1/r\bigr)d(h,h')
 \le Cr\epsilon^{-6}d(h,h').
 \label{ch:later-fisher}
\end{equation}

It remains to check the regularity promised in the channel definition.
For a fixed finite input support $S$, only finitely many centers $z$ can
have $D_z\ne0$, all in the $2r$ neighborhood of $S$.  For each such
center, the alphabet and string set are finite and independent of $P$.
The functions $S(P;z,w)$ and $\pi_z(P)_w$ are smooth functions of the
coordinates of $P$: the scores are linear, and each exponential sum has a
strictly positive denominator.  Thus \eqref{ch:later-channel} is a fixed
finite smooth formula on each enlarged support space.  If $L_{z,w}=0$,
strict positivity of $\pi_z(w)$ implies $D_z=0$; also $\dot D_z=0$ by
Lemma~\ref{ch:selector-estimates}.  Consequently $\dot L_{z,w}=0$.
This proves smoothness at faces and the required behavior at zero entries.

\begin{proof}[Completion of the proof of Proposition~\ref{ch:summary}]
Use the same selector $A$ in both versions.  Inequality
\eqref{ch:selector-concentration} gives
\eqref{ch:selection-bound}.  The early construction
\eqref{ch:early-channel} gives the two prediction inequalities with
errors $0$ and $2\epsilon$; the later construction
\eqref{ch:later-channel} gives errors $5\epsilon$ and $6\epsilon$.
All are bounded by $10\epsilon$.
The Fisher estimates \eqref{ch:selector-fisher},
\eqref{ch:early-fisher}, and \eqref{ch:later-fisher} imply
\eqref{ch:summary-fisher}, since $r\ge1$ and $\epsilon\le1$.
The preceding constructions also verify the channel and cross-kernel
requirements, including the assertions at zero coordinates.
\end{proof}

For use with differential forms, regard the labels of $A$ and $L$ as
disjoint even when their underlying sets happen to coincide.  On their
disjoint union define the symmetric kernel $K$ by
\begin{equation}
 \begin{split}
 K(y,\omega)=K(\omega,y)&=\tfrac12\kappa(y,\omega)
       \quad(y\in G,\ \omega\in\Lambda),\\
 K&=0\quad\text{within either channel}.
 \end{split}
 \label{ch:symmetric-kernel}
\end{equation}
For vectors on this union, write $K(U,V)=\sum_{\xi,\eta}K(\xi,\eta)U_\xi V_\eta$.
The combined field $Y=A\oplus L$ has mass at most two, and
\begin{equation}
 K(Y(P),Y(P))
    =\sum_yA(P)_y\operatorname{pred}(P,y).
 \label{ch:kernel-diagonal}
\end{equation}
The kernel is bounded and invariant under simultaneous left translation.
The expression in \eqref{ch:fisher-definition} also makes sense for this
field of mass at most two, and its square is
$\mathcal F_A(P;h,h')^2+\mathcal F_L(P;h,h')^2$.
Thus Proposition~\ref{ch:summary} controls its derivatives as well.
These facts permit both channels to enter a single bilinear pairing in the
averaging construction.
